\section{Introduction}
\IEEEPARstart{T}{he} advancement of Industry 4.0 has introduced a new paradigm in the manufacturing sector through the seamless integration of physical and digital systems \cite{soori2024}. A key concept emerging from this transformation is virtual manufacturing, a production approach that integrates sensor data, digital models, and intelligent devices to achieve real-time process optimization \cite{soori2024}. Within the context of modern manufacturing, virtual inspection and quality control systems play a critical role\cite{soori2024}. These systems leverage imaging technology and visual data processing to monitor, analyze, and guarantee product quality without direct physical intervention \cite{soori2024}. In material extrusion systems like 3D printers, apart from chemical and thermal extrusion parameters, the absolute positioning accuracy of the machine's mechanical movements significantly influences the quality of the printed model \cite{zaborniak2024}. With prolonged utilization, the pulley and belt linear motion transmission systems can undergo calibration drift due to structural relaxation and mechanical wear \cite{xu2020, zaborniak2024, papageorgiou2019}. Traditional calibration methods for high-precision machinery employ a mechanical dial gauge for error profiling by evaluating the deviation between the actual displacement and the commanded position \cite{iso230_1_2012}. Despite offering a fine measurement resolution down to 10~$\mu$m, this manual approach suffers from several drawbacks, including susceptibility to human error and poor time efficiency.

Several prior studies have developed automated calibration methods using physical sensors for 3D printers; however, such hardware automation incurs a significant cost increase. Examples of research attempting to automate error profiling and 3D printer calibration include inductive methods \cite{Schouten2022}, camera-based systems using Hue Saturation Value (HSV) color segmentation algorithms \cite{Chen2025}, laser triangulation \cite{Munteanu2021}, and computer vision integrated with planar ArUco markers \cite{Owais2026}. From an implementation cost perspective, the computer vision approach combined with ArUco represents the most cost-effective solution, as it only requires a standard camera and paper-printed markers. Nevertheless, this method falls short in terms of accuracy, exhibiting a measurement tolerance error of 84~$\mu$m. This accuracy level is inferior to that of the inductive method, which achieves a translational accuracy of 45~$\mu$m with a repeatability of 2~$\mu$m \cite{Schouten2022}. Hence, there is an urgent need to develop an automated kinematic calibration approach for desktop 3D printers that is highly accurate, precise, and cost-effective.

To address this need, this study proposes the concept of a virtual encoder system that adopts the fundamental principles of marker-based position tracking, leveraging computer vision and fiducial markers. The virtual encoder is designed to translate camera pixel shifts into actual machine displacements, effectively substituting the function of a physical rotary encoder. By employing this system, machine parameters such as steps per millimeter, alongside geometric and kinematic errors like backlash and skewness, can be evaluated by analyzing the positional deviations along the machine axes captured by the camera. The resulting error profiling data is subsequently utilized to execute error compensation during machine calibration. This study focuses primarily on XY-axis calibration, as Z-axis alignment can already be reliably achieved using electromechanical limit switches or inductive proximity sensors \cite{Schouten2022}.